% !TeX program = xelatex
% ============================================================
%  StudyMind——基于大模型的个性化 AI 学习助手 · 应用方案文档
%  编译：TeXstudio 选择 XeLaTeX（文件头已自动指定），F5 编译查看
%  截图目录：D:\ai学习助手\（已通过 \graphicspath 指定）
% ============================================================
\documentclass[11pt,a4paper]{ctexart}

\usepackage[top=2.5cm,bottom=2.5cm,left=2.8cm,right=2.8cm]{geometry}
\usepackage{amsmath,amssymb}
\usepackage{booktabs}
\usepackage{array}
\usepackage{xcolor}
\usepackage{enumitem}
\usepackage{graphicx}
\usepackage{caption}
\usepackage{fancyhdr}
\usepackage[hidelinks]{hyperref}

% ---- 图片目录（截图都在 D:\ai学习助手） ----
\graphicspath{{D:/ai学习助手/}}

% ---- 手动编号的章节（每个一级章节自动另起一页；目录显示干净） ----
\newcommand{\SX}[1]{\clearpage\phantomsection\section*{#1}\addcontentsline{toc}{section}{#1}}
\newcommand{\SSX}[1]{\phantomsection\subsection*{#1}\addcontentsline{toc}{subsection}{#1}}

% ---- 插图 + 题注 ----
\captionsetup{font=small,labelsep=quad}
\newcommand{\fig}[2]{%
  \begin{center}
    \includegraphics[width=0.82\textwidth,height=0.55\textheight,keepaspectratio]{#1}
    \captionof{figure}{#2}
  \end{center}
}

% ---- 目录只到二级 ----
\setcounter{tocdepth}{2}

% ---- 页眉页脚 ----
\pagestyle{fancy}
\fancyhf{}
\fancyhead[L]{\small StudyMind —— 基于大模型的个性化 AI 学习助手}
\fancyhead[R]{\small 应用方案文档}
\fancyfoot[C]{\small 第 \thepage 页}
\renewcommand{\headrulewidth}{0.4pt}

\setlist{nosep}

\begin{document}

% ==================== 封面 ====================
\begin{titlepage}
  \centering
  \vspace*{2.5cm}
  {\color{black!70}\large 23 级人工智能专业课程设计 · 应用方案文档}\\[1.8cm]
  {\Huge\bfseries StudyMind}\\[0.6cm]
  {\LARGE\bfseries ——基于大模型的个性化 AI 学习助手——}\\[1.2cm]
  {\large 核心创新：基于知识图谱的个性化学习路径动态调整（GraphRAG）}\\[3.2cm]
  {\large 课程：《云计算与大数据实践》}\\[0.5cm]
  {\large 班级：23 级人工智能 1 班}\\[0.5cm]
  {\large 成员：朱晨馨 \quad 游星宇}\\[1.5cm]
  {\large 2026 年 9 月}
\end{titlepage}

% ==================== 目录 ====================
\tableofcontents
\newpage

% ==================== 一、项目背景与痛点分析 ====================
\SX{一、项目背景与痛点分析}

\SSX{1.1 项目背景}

这几年在线教育发展很快，大学生手里的学习资料越来越多：电子教材、课件 PPT、网课视频、各种讲义，但资料多并没有让学习变轻松，复习的时候仍然不知道该从哪里下手。另一方面，大语言模型（LLM）近两年的进步非常快，文本理解、推理和结构化输出的能力已经比较可靠，调用成本也在下降，这让“用 AI 做个性化学习”从想法变成了可以落地的工程问题。

StudyMind 就是围绕这个想法做的一个系统：学生把教材交给系统，系统自动解析教材、生成知识图谱和学习计划、按掌握度出题，答错之后还会分析错因并调整后续的学习任务。整个系统围绕“解析—建图—规划—练习—诊断—调整”这条主线设计。

\SSX{1.2 痛点问题}

\textbf{（一）资料多而散，形不成知识体系。} 教材、PPT、网课各有各的组织方式，复习时东翻一点西翻一点，经常学完一章也说不清它在整个课程里的位置，更谈不上把握知识点之间的前后关系。

\textbf{（二）刷题盲目，找不到薄弱点。} 大量练习之后还是说不清自己哪里弱、为什么弱；错题基本靠手工整理，整理的也只是一道道孤立的题，很少有人去分析错误背后的原因，同样的错误反复出现。

\textbf{（三）学习计划是固定的。} 课表和大多数学习 App 的计划都是统一顺序、统一进度，不会因为个人的掌握程度改变。计划一旦生成就按部就班执行，跟“因材施教”差得很远。

这三个问题归结起来是一件事：现有工具负责搬运内容，但不负责“教”。StudyMind 的设计目标，就是用知识图谱解决知识体系问题，用 Elo 掌握度模型解决薄弱点定位问题，用 GraphRAG 诊断解决学习路径调整问题。

% ==================== 二、需求分析 ====================
\SX{二、需求分析}

\SSX{2.1 目标用户群体}

\begin{center}
\small
\renewcommand{\arraystretch}{1.3}
\begin{tabular}{p{2.6cm}p{5.4cm}p{6.0cm}}
  \toprule
  \textbf{用户群体} & \textbf{典型场景} & \textbf{核心诉求} \\
  \midrule
  高校在校大学生 & 理工科专业课预习、期末复习 & 快速建立课程知识体系、考前精准查漏补缺 \\
  考研/证书备考学生 & 专业课长周期备考 & 长线计划管理、薄弱点动态定位、错题闭环 \\
  职业培训/终身学习者 & 技能类知识自学 & 零散资料体系化、碎片时间高效练习 \\
  \bottomrule
\end{tabular}
\end{center}

\SSX{2.2 功能需求}

\begin{center}
\small
\renewcommand{\arraystretch}{1.3}
\begin{tabular}{p{1.6cm}p{3.4cm}p{9.0cm}}
  \toprule
  \textbf{类别} & \textbf{功能模块} & \textbf{功能说明} \\
  \midrule
  基础功能 & 教材上传与智能解析 & 支持 PDF 上传（含扫描版 OCR），自动提取章节结构与知识点 \\
  基础功能 & 无教材模式 & 用户没有 PDF 时，输入学科主题即可由 LLM 生成知识骨架 \\
  基础功能 & 知识图谱生成与可视化 & 自动抽取知识点及前置/关联关系，交互式图谱浏览 \\
  基础功能 & 学习计划生成 & 双引擎计划：拓扑排序定结构，LLM 写文案，按天拆解任务 \\
  基础功能 & 智能答题练习 & 每日任务/知识点专项/阶段总考三种模式；客观题规则判分、主观题 LLM 按评分细则批改 \\
  基础功能 & 掌握度看板 & Elo 掌握度、连对统计、薄弱点排行可视化展示 \\
  基础功能 & 激励统计栏 & 连续打卡、今日学习时长、目标进度等激励数据全页面实时联动 \\
  创新功能 & \textbf{个性化学习路径动态调整（GraphRAG）} & 答错自动触发诊断：图谱邻域检索、掌握度状态、LLM 根因分析，自动插入补学/复习任务，形成“答错—诊断—重排—再练—掌握”闭环 \\
  创新功能 & Elo 自适应出题 & 根据掌握度与最新诊断报告个性化调整题目难度、认知层次与干扰项 \\
  创新功能 & AI 助教问答 & 围绕题目与知识点进行启发式答疑，支持一键跳转知识图谱 \\
  \bottomrule
\end{tabular}
\end{center}

\SSX{2.3 非功能需求}

除了功能本身，系统还有几条工程上的要求：

\textbf{（1）响应速度。} 提交答案的接口要在 1 秒内返回判分结果（当前实测在 0.02 秒左右）；出题、诊断这类要调大模型的操作都放到后台任务里执行，前端轮询取结果，不阻塞页面。

\textbf{（2）稳定可用。} 后台任务要能看到进度；加载失败可以重试，页面不白屏；答题流程做成状态机，提交按钮在异常情况下也能复位，不会一直转圈。

\textbf{（3）数据安全。} 题目答案只保存在服务端，答题接口不返回答案；API 密钥放在 \texttt{.env} 里，不进代码库。

\textbf{（4）成本可控。} 同一个用户答过的题直接复用题库，题干做向量去重，预生成题目限量，这三条都是为了控制大模型的调用费用。

\textbf{（5）便于扩展。} 后端按模块组织，任务管理器目前是内存实现，接口与 Celery 一致，以后可以替换；三个数据库各管各的，可以独立扩容。

% ==================== 三、开发工具与技术方案 ====================
\SX{三、开发工具与技术方案}

\SSX{3.1 开发工具与环境}

\begin{center}
\small
\renewcommand{\arraystretch}{1.3}
\begin{tabular}{p{2.0cm}p{5.4cm}p{6.6cm}}
  \toprule
  \textbf{层次} & \textbf{选型} & \textbf{说明} \\
  \midrule
  开发工具 & VS Code、Git、Docker Compose & 统一开发环境，容器化部署 \\
  前端框架 & Vue 3 + TypeScript + Vite & 组件化开发、类型安全 \\
  前端 UI & Tailwind CSS + Element Plus & 原子化样式 + 成熟组件库 \\
  前端可视化 & AntV G6、ECharts & 知识图谱渲染、掌握度图表 \\
  状态管理 & Pinia & 全局激励统计等跨页共享状态 \\
  后端框架 & Python 3 + FastAPI + Uvicorn & 异步 Web 框架，自动生成接口文档 \\
  LLM 编排 & LangChain & 结构化输出（function\_calling）与提示词管理 \\
  关系数据库 & PostgreSQL & 业务事实源（用户/教材/题目/答题记录/掌握度） \\
  图数据库 & Neo4j & 知识图谱存储与邻域查询（GraphRAG 检索） \\
  向量数据库 & ChromaDB & 教材 chunk 向量索引（RAG 检索） \\
  嵌入模型 & BGE-M3（本地部署） & 中英双语向量化，离线可用、零调用成本 \\
  大模型 & DeepSeek（OpenAI 兼容协议） & 知识抽取/出题/诊断/批改等生成任务 \\
  文档解析 & PyMuPDF + pdfplumber + pypdf + Tesseract OCR & 多策略解析链，兼容扫描版教材 \\
  \bottomrule
\end{tabular}
\end{center}

\SSX{3.2 系统架构设计}

系统分成四层：

\textbf{（1）表现层。} Vue3 单页应用，共五个页面：知识库管理、知识图谱、学习计划、答题练习、掌握度看板。前端通过 RESTful API 和后端交互，异步任务统一轮询 \texttt{/api/task\_status/\{id\}} 获取进度。

\textbf{（2）业务逻辑层。} FastAPI 接口层加服务层（教材解析、知识构建、计划生成、出题判分、诊断引擎），还有一个轻量任务管理器：内存注册表加后台线程。所有要调大模型的操作都放进后台任务执行，接口立刻返回 \texttt{task\_id}，前端拿这个 id 轮询结果。

\textbf{（3）数据层。} 三个数据库分开用：PostgreSQL 存业务事实（教材、知识点、题目、答题记录、掌握度状态）；Neo4j 存知识图谱，负责邻域查询和可视化数据；ChromaDB 存教材切片的向量，给 RAG 检索用。

\textbf{（4）AI 服务层。} DeepSeek 大模型负责知识抽取、出题、诊断、批改这些生成任务（统一封装成结构化输出）；BGE-M3 嵌入模型本地部署，负责向量化；Tesseract OCR 兜底扫描版教材。

\fig{系统整体架构图.png}{系统整体架构图}

\SSX{3.3 核心技术方案}

\subsubsection*{教材解析与 RAG 检索增强生成}

教材上传后先走一条多策略解析链：优先用 PyMuPDF 提取文本，失败依次降级到 pdfplumber 和 pypdf；如果教材是扫描版、没有文本层，就转 Tesseract OCR（中文 \texttt{chi\_sim}，150dpi），实测 300 多页的扫描教材可以完整解析。解析时优先利用 PDF 书签组织章节。解析出来的内容按语义切成 chunk，用本地部署的 BGE-M3 算向量后存入 ChromaDB。出题、诊断、答疑这些场景需要教材原文时，先从向量库召回相关 chunk 放进提示词，让大模型的输出有教材依据，减少胡编乱造。

\subsubsection*{知识图谱构建与可视化}

LLM 对教材内容做结构化抽取，输出知识点（名称、定义）和知识点之间的关系（前置、关联）。没有教材的时候也可以直接让 LLM 按学科主题生成一个知识骨架（无教材模式）。抽取结果同时写进 PostgreSQL 和 Neo4j。前端用 AntV G6 渲染图谱，支持缩放、拖拽和节点点击。

\subsubsection*{基于 Elo 模型的掌握度评估与动态路径调整}

\textbf{（1）Elo 掌握度模型。} 每个“用户—知识点”对维护一个掌握度分数 $\theta$，初始 1500。每答一题按 Elo 公式更新：答对概率
\[
P = \frac{1}{1 + 10^{-(\theta - \beta)/400}}
\]
其中 $\beta$ 是题目难度映射过来的值；答对则 $\theta \leftarrow \theta + K(1-P)$，答错则 $\theta \leftarrow \theta - K \cdot P$，$K$ 取 32。这里有个容易踩的坑：题目难度是 1--5，$\theta$ 是 1500 量级，直接套公式会出现“答对也不涨分”的问题，所以做了一个 \texttt{difficulty\_to\_elo} 映射把难度换算到 Elo 量纲。掌握度还有连对约束：连续答对 $\geq 3$ 题或 $\theta$ 达到 1700 就算“已掌握”，计划里该知识点的剩余任务自动跳过。

\textbf{（2）GraphRAG 错题诊断。} 学生答错后，系统在后台启动诊断任务：以错题知识点为中心，从 Neo4j 里取它的图谱邻域，再结合这个学生全部知识点的掌握度分布，交给 LLM 分析错误根因，输出补学路径（前置知识点链条）和易混淆概念，存成诊断报告。

\textbf{（3）计划动态重排。} 诊断报告驱动调整引擎改计划：补学路径里的知识点不在计划中就插入学习任务，排期太靠后就把任务提前；易混淆概念插入对比复习任务；插入时按每日任务上限均衡，避免某一天任务堆满。

\textbf{（4）自适应出题。} 下次给这个知识点出题时，读取最新的诊断报告，按报告里的策略调整题目难度、认知层次（Bloom）和干扰项。这样整个链条就闭合了：答错—诊断—重排—降难度—再练—掌握。这是本项目和一般题库类产品最大的区别。

% ==================== 四、作品功能说明与 AI 核心作用 ====================
\SX{四、作品功能说明与 AI 核心作用}

\SSX{4.1 教材上传与智能解析}

首页可以上传 PDF 教材，也可以不传教材、直接输入学科主题让 AI 生成知识骨架（无教材模式）。上传后解析在后台执行，页面用进度条和阶段提示反馈，还会根据进度变化预估剩余时间。

AI 在这个环节做的事：

\textbf{（1）解析。} 普通电子版教材用 PyMuPDF 主链路解析，遇到异常自动降级 pdfplumber、pypdf；扫描版教材转 Tesseract OCR（\texttt{chi\_sim}，150dpi），并优先用 PDF 书签组织章节结构，实测 300 多页扫描教材能完整解析。

\textbf{（2）切块与向量化。} 章节内容按语义切 chunk，本地 BGE-M3 嵌入后存 ChromaDB，后面出题、诊断、答疑都靠它做检索。

\textbf{（3）知识点抽取。} LLM 从教材内容里识别知识点实体（名称加定义），用 \texttt{function\_calling} 约束输出格式，去重后双写 PostgreSQL 和 Neo4j。

\fig{教材上传与解析界面截图.png}{教材上传与解析界面}

\SSX{4.2 知识图谱生成}

系统自动构建“知识点 + 前置/关联关系”的知识图谱，图谱页基于 AntV G6 渲染，支持缩放、拖拽、节点聚焦，薄弱知识点标红。

AI 做的事：

\textbf{（1）关系推理。} LLM 抽取知识点时同步判断关系：前置关系（学 B 前必须掌握 A）和关联关系（概念相近、容易混淆），输出结构化三元组写入 Neo4j。关键词抽取工具做不到这种语义判断。

\textbf{（2）图谱服务于诊断。} 图谱不是摆设，错题诊断（GraphRAG）要靠它：学生答错时，系统沿图谱的边找到“该补哪些前置知识点”。

\textbf{（3）可视化。} 点开节点能看到知识点定义，薄弱点分布一眼可见。

\fig{知识图谱可视化截图.png}{知识图谱可视化页面}

\SSX{4.3 学习计划与智能出题}

学习计划由两部分拼起来：结构由算法定，文案由 LLM 写。算法部分对知识图谱做拓扑排序，保证前置知识点一定排在后继之前，不会出现“没学前置就学后继”；LLM 负责给每天的任务写说明文案。用户可以设定总天数、每日时长和每日任务上限，计划按天拆分，跟踪每个任务的完成状态。

出题方面：

\textbf{（1）复用优先。} 该用户没答过的题库题直接复用，不够才调 LLM 生成（3 路并发），生成时把已出过的题干注入提示词、再做一轮向量去重，防止题目雷同。

\textbf{（2）按诊断出题。} 出题前读该知识点最新的诊断报告，按报告调整难度、认知层次（Bloom）和干扰项。

\textbf{（3）判分在服务端。} 支持单选、多选、判断、填空、简答；客观题服务端规则判分，简答题让 LLM 按评分细则逐点批改并写评语；答案不出接口，防止在浏览器里直接看到答案。

\fig{学习计划与答题界面截图.png}{学习计划与答题练习界面}

\SSX{4.4 掌握度看板}

看板按知识点展示 Elo 掌握度、连对数、答对答错次数和最近复习时间，按掌握度排序，薄弱点标红。顶部的激励栏（连续打卡、今日学习时长、目标进度）和看板数据全局联动，答完题立刻刷新。

值得说明的是，看板上的掌握度不是简单数对错：每题判分后 Elo 即时更新（接口亚秒返回），$\theta$ 的变化每题都有记录，能看出具体是在哪道题、哪个知识点上涨分的。

\fig{掌握度看板截图.png}{掌握度看板}

\SSX{4.5 创新功能：个性化学习路径动态调整（GraphRAG）}

这是本项目的核心创新点。一般题库类产品做到“判对错、给解析”就结束了，学习计划是死的。StudyMind 往前多走了几步：AI 不仅要判断对错，还要分析错因、给出补学方案，并且直接改学习计划。

完整流程：

\textbf{（1）答错即触发，不阻塞。} 判分接口亚秒返回对错和 Elo 变化，同时后台启动诊断任务，学生继续答题。

\textbf{（2）GraphRAG 检索。} 以错题知识点为中心取 Neo4j 邻域（前置/关联节点），加上该生全部知识点的掌握度分布（$\theta$ 图谱），拼成诊断上下文。

\textbf{（3）LLM 根因诊断。} 大模型输出结构化诊断报告：错误根因、补学路径、易混淆概念、下次出题策略。

\textbf{（4）计划自动重排。} 补学路径里的知识点不在计划则插入学习任务，排期靠后则提前；易混淆概念插入对比复习任务；按每日任务上限均衡插入。

\textbf{（5）再练直到掌握。} 下次出题按诊断策略降难度、调认知层次；$\theta$ 到 1700 或连对 3 题算已掌握，剩余任务自动跳过。

举个例子：学生在“OSI 参考模型”上答错，$\theta$ 从 1468.5 降到 1445.3；AI 沿图谱诊断出根因是前置知识点“网络分层思想”没掌握，自动在当天计划里插入补学任务和对比复习题；补学之后再练，$\theta$ 回升达标，系统跳过后续重复任务。

\fig{动态调整错题解析界面截图.png}{答错后 AI 诊断与计划动态调整}

AI 大模型在五个模块里分别起的作用：

\begin{center}
\small
\renewcommand{\arraystretch}{1.3}
\begin{tabular}{p{3.0cm}p{11.0cm}}
  \toprule
  \textbf{模块} & \textbf{AI 大模型的作用} \\
  \midrule
  教材解析 & 章节组织、知识点结构化抽取、扫描件 OCR 兜底 \\
  知识图谱 & 前置/关联关系推理、语义结构化三元组 \\
  计划与出题 & 计划文案生成、自适应出题（难度/Bloom/干扰项）、主观题批改 \\
  掌握度看板 & 基于 Elo 的连续掌握度状态建模与实时更新 \\
  动态调整 & GraphRAG 根因诊断、补学路径推理、计划自动重排 \\
  \bottomrule
\end{tabular}
\end{center}

% ==================== 五、使用说明与交互指南 ====================
\SX{五、使用说明与交互指南}

\SSX{5.1 新用户入门流程}

\textbf{第 1 步：注册/登录。} 打开 StudyMind，注册并登录账号（课程设计演示版本内置演示账号，可以直接登录）。

\textbf{第 2 步：创建知识库。} 进入首页，二选一：方式 A，点击“上传教材”，选择一个 PDF 文件；方式 B，无教材模式，输入学科和主题（比如“数据结构 · 二叉树”），AI 直接生成知识骨架。

\textbf{第 3 步：等待 AI 解析。} 上传后页面显示解析进度和预计剩余时间，期间可以浏览其他页面；解析完成后自动跳转知识图谱。

\textbf{第 4 步：浏览知识图谱。} 查看自动生成的知识点和前置/关联关系，确认知识库构建成功。

\SSX{5.2 核心功能操作指南}

\textbf{（1）如何上传教材并生成图谱。}

进入首页，点“上传教材”，选择 PDF，等待解析进度完成，点“查看图谱”，缩放、拖拽浏览知识体系，点任意节点查看该知识点的定义和关联。

\fig{上传解析与图谱浏览操作截图.png}{教材上传与图谱浏览}

\textbf{（2）如何制定并开始每日学习计划。}

进入“学习计划”页，选择教材，设置总天数（比如 7 天）、每日学习时长（比如 45 分钟）和每日任务上限，点“生成计划”，AI 生成按天拆解的学习计划，点“去学习”进入答题练习。计划生成期间能看到进度提示；生成后如果时间安排有变，可以一键“重新均衡”任务分布。

\textbf{（3）如何进行答题练习并查看 AI 解析。}

进入“答题练习”页，选一种模式：每日任务（完成当天计划里的题）、知识点专项（从薄弱点优先列表里挑一个知识点集中强化）、阶段总考（覆盖整本教材的模拟测试，带 45 分钟倒计时）。

答题流程：读题，点选项，点“提交本题”，马上看到对错、Elo 变化和分层解析（先给一句概要，可展开看详解）；答错时 AI 自动生成学习诊断（约 10--30 秒，不影响继续答题）；点“下一题”继续；最后一题点“查看总结”，看正确率、掌握度变化和错题回顾。顶部激励栏的连续答对、今日学习时长随作答实时刷新。

\fig{答题练习与 AI 解析操作截图.png}{答题练习与 AI 解析}

\textbf{（4）如何查看看板并根据 AI 建议调整复习。}

进入“掌握度看板”，按 $\theta$ 值看知识点掌握度排序，关注标红的薄弱点。复习有两种方式：直接点薄弱点发起“专项练习”；或者回到答题页，按 AI 诊断给出的补学路径依次复习前置知识点。系统会自动把诊断出的补学任务插进每日计划，跟着更新后的计划走就行。

\fig{掌握度看板操作截图.png}{掌握度看板与薄弱点复习}

% ==================== 六、应用前景与商业模式 ====================
\SX{六、应用前景与商业模式}

\SSX{6.1 应用前景}

\textbf{（1）高校教育。} 高校课程教材统一、知识点成体系，是最好落地的场景：学生期末复习、教师布置练习、补考重修辅导都能用上。系统沉淀的班级级掌握度数据对教师也有用，能看出全班普遍薄弱的点在哪里。

\textbf{（2）职业培训。} IT 技能、资格考试、企业内训这些场景知识更新快、学员水平参差，最需要个性化培训。培训机构的教研人员可以省下大量出题和批改的时间。

\textbf{（3）终身学习。} 个人自学时“学什么、怎么学、学到什么程度”是最常见的问题，StudyMind 的知识体系构建和动态路径调整能力可以直接用上，有潜力做成通用型的 AI 学习助手。

\SSX{6.2 商业模式}

\textbf{（1）Freemium 订阅。} 基础功能（教材解析、图谱、基础练习、看板）免费，高级 AI 功能（GraphRAG 深度诊断、更多出题额度、简答题 AI 批改、长周期计划）订阅收费。大模型调用成本高的功能放付费档，保证成本模型成立。

\textbf{（2）B 端授权。} 面向高校院系做课程级部署：教师上传教材，全班共用图谱和题库，教师端有班级掌握度总览；面向教辅出版机构，把教材内容做成交互式学习产品，按课程或学期授权收费。

\textbf{（3）增值服务。} 面向备考用户提供周期性的学习诊断报告（薄弱点分析、掌握度趋势、复习建议）；面向家长和教师提供学情报告；后续可以接 AI 助教按次计费、人工答疑等服务。

% ==================== 结语 ====================
\SX{结语}

StudyMind 把“教材解析—知识图谱—学习计划—答题练习—诊断调整”整条流程做了出来，其中个性化学习路径动态调整这个创新点已经可以完整演示。项目里用到了 RAG、知识图谱、大模型结构化输出、Elo 模型这些课程核心知识，前端、后端和数据库的工程实现也比较完整。接下来还有很多可以完善的地方，比如真正的用户体系、任务队列持久化、更精细的难度模型，希望在后续迭代里继续做下去。

\end{document}
